\documentclass[12pt]{article}
\usepackage[utf8]{inputenc}
\usepackage[T1]{fontenc}
\usepackage[margin=1in]{geometry}
\usepackage{setspace}
\usepackage{xcolor}
\usepackage[colorlinks=true, urlcolor=blue!60!black, linkcolor=blue!60!black]{hyperref}
\setlength{\parindent}{2em}
\setlength{\parskip}{1.25em}
\renewcommand{\baselinestretch}{1.0}
\begin{document}
\begin{center}
{\LARGE \textit{Hammer v. Dagenhart}, 247 U.S. 251 (1918)\par}
\vspace{0.8em}
{\large Case Brief\par}
\end{center}
\vspace{1em}

\section{Memory Jogger}
Congress exceeded its Commerce Clause power and invaded the Tenth Amendment police powers reserved to the States when it barred the interstate shipment of harmless goods produced in factories employing child labor---a dual-federalism holding later expressly overruled in \textit{United States v. Darby} (1941).

\section{Detailed Case Facts}
On September 1, 1916, Congress enacted the Keating-Owen Child Labor Act, ch.\ 432, 39 Stat.\ 675, which prohibited any producer, manufacturer, or dealer from shipping or delivering for shipment in interstate or foreign commerce any product of a mine or quarry in which children under age sixteen had been employed within thirty days prior to removal, or any product of a mill, cannery, workshop, factory, or manufacturing establishment in which, within thirty days prior to removal, children under age fourteen had been employed or children between ages fourteen and sixteen had worked more than eight hours in a day, more than six days in a week, or between 7:00 p.m.\ and 6:00 a.m.\ (\href{https://supreme.justia.com/cases/federal/us/247/251/}{247 U.S. at 268--69, 277 n.*}; \href{https://www.courtlistener.com/opinion/99190/hammer-v-dagenhart/}{CourtListener}). At the time, every state had enacted some form of child labor restriction, and North Carolina law prohibited the employment of children under age twelve (\href{https://supreme.justia.com/cases/federal/us/247/251/}{247 U.S. at 275}).

Roland Dagenhart was a cotton-mill worker in Charlotte, North Carolina, whose two minor sons---one under age fourteen and one between ages fourteen and sixteen (Reuben Dagenhart)---worked alongside him in a cotton mill (\href{https://supreme.justia.com/cases/federal/us/247/251/}{247 U.S. at 268}; \href{https://fedsoc.org/case/hammer-v-dagenhart}{Federalist Society}). Because the Keating-Owen Act barred the interstate shipment of goods from any mill employing the younger son at all or the older son for the mill's regular hours, enforcement of the federal statute threatened their continued employment (\href{https://www.abajournal.com/magazine/article/june_3_1918_child_labor_law_declared_unconstitutional}{ABA Journal}). Acting in his own behalf and as next friend of his two minor sons, Roland Dagenhart sued W. C. Hammer, the United States Attorney for the Western District of North Carolina, seeking an injunction against enforcement of the Act on three constitutional grounds: that it was not a valid regulation of interstate or foreign commerce under Article I, Section 8, Clause 3; that it contravened the Tenth Amendment; and that it violated the Fifth Amendment (\href{https://supreme.justia.com/cases/federal/us/247/251/}{247 U.S. at 268--69}).

\section{Procedural History}
Dagenhart filed a bill in equity in the United States District Court for the Western District of North Carolina to enjoin U.S.\ Attorney W. C. Hammer from enforcing the Keating-Owen Act (\href{https://supreme.justia.com/cases/federal/us/247/251/}{247 U.S. at 268}). The District Court held the Act unconstitutional and entered a decree enjoining its enforcement (\href{https://supreme.justia.com/cases/federal/us/247/251/}{247 U.S. at 268}). Hammer took a direct appeal to the United States Supreme Court (\href{https://supreme.justia.com/cases/federal/us/247/251/}{247 U.S. at 251, 268}). The case was argued on April 15 and 16, 1918, and decided on June 3, 1918, with the Supreme Court affirming the District Court's decree by a 5--4 vote (\href{https://supreme.justia.com/cases/federal/us/247/251/}{247 U.S. at 251, 277}; \href{https://www.oyez.org/cases/1900-1940/247us251}{Oyez}).

\section{Judicial Votes}
\textbf{Vote:} 5--4, affirming the District Court's decree invalidating the Keating-Owen Child Labor Act (\href{https://www.oyez.org/cases/1900-1940/247us251}{Oyez}; \href{https://supreme.justia.com/cases/federal/us/247/251/}{247 U.S. at 268, 277, 281}).
\begin{itemize}
\item \textbf{Majority (5):} Justice William R. Day (author), joined by Chief Justice Edward D. White, Justice Willis Van Devanter, Justice Mahlon Pitney, and Justice James C. McReynolds (\href{https://www.oyez.org/cases/1900-1940/247us251}{Oyez}; \href{https://supreme.justia.com/cases/federal/us/247/251/}{247 U.S. at 268}).
\item \textbf{Concurrences (0):} None.
\item \textbf{Dissent (4):} Justice Oliver Wendell Holmes, Jr.\ (author), joined by Justice Joseph McKenna, Justice Louis D. Brandeis, and Justice John H. Clarke (\href{https://supreme.justia.com/cases/federal/us/247/251/}{247 U.S. at 277, 281}).
\end{itemize}

\section{Holding}
\textbf{Judgment:} Affirmed. The Keating-Owen Child Labor Act of 1916 is unconstitutional (\href{https://supreme.justia.com/cases/federal/us/247/251/}{247 U.S. at 277}).

\textbf{Rule Announced:} Congress lacks authority under the Commerce Clause to prohibit the interstate transportation of ordinary, intrinsically harmless manufactured goods merely because they were produced under local labor conditions (such as child labor) that Congress disapproves. Manufacturing is a local activity distinct from interstate commerce, and Congress may prohibit the interstate movement of goods only when the goods themselves are inherently harmful or interstate transportation is necessary to accomplish the harmful result (\href{https://supreme.justia.com/cases/federal/us/247/251/}{247 U.S. at 270--74, 276}). Furthermore, using the commerce power to standardize local labor conditions or equalize economic competition among states invades the police power reserved to the States under the Tenth Amendment (\href{https://supreme.justia.com/cases/federal/us/247/251/}{247 U.S. at 273--76}).

\textbf{Holding vs.\ Dicta:} The narrow holding is that Congress may not use the Commerce Clause to bar the interstate shipment of harmless manufactured commodities produced with child labor in order to regulate local manufacturing hours and ages. Justice Day's broader warnings---that sustaining such a law would mean ``all freedom of commerce will be at an end'' and ``our system of government be practically destroyed'' (\href{https://supreme.justia.com/cases/federal/us/247/251/}{247 U.S. at 276}), as well as his assertion that the Tenth Amendment reserves to the States all powers not ``expressly'' delegated to the National Government (\href{https://supreme.justia.com/cases/federal/us/247/251/}{247 U.S. at 275}) and his hypothetical about minimum-wage and maximum-hour laws for women (\href{https://supreme.justia.com/cases/federal/us/247/251/}{247 U.S. at 273})---constitute dicta. (Notably, the Supreme Court unanimously overruled \textit{Hammer v. Dagenhart} twenty-three years later in \href{https://supreme.justia.com/cases/federal/us/312/100/}{\textit{United States v. Darby}, 312 U.S. 100, 116--17 (1941)}.)

\section{Analysis of Opinions}

\subsection{Majority Opinion (Justice Day)}
Justice Day began by invoking Chief Justice Marshall's definition in \textit{Gibbons v. Ogden}, 22 U.S.\ (9 Wheat.) 1 (1824), that the commerce power is the power ``to prescribe the rule by which commerce is to be governed''---which Day paraphrased as the power to control the means by which commerce is carried on, rather than a general power to prohibit and destroy commerce in ordinary commodities (\href{https://supreme.justia.com/cases/federal/us/247/251/}{247 U.S. at 269--70}). He acknowledged prior decisions sustaining federal prohibitions on interstate movement---including \textit{Champion v. Ames} (lottery tickets), \textit{Hipolite Egg Co. v. United States} (adulterated food and drugs), \textit{Hoke v. United States} and \textit{Caminetti v. United States} (interstate transportation of women for prostitution or debauchery), and \textit{Clark Distilling Co. v. Western Maryland Ry. Co.} (intoxicating liquor)---but distinguished them on a single doctrinal axis: in each of those cases, ``the use of interstate transportation was necessary to the accomplishment of harmful results'' (\href{https://supreme.justia.com/cases/federal/us/247/251/}{247 U.S. at 270--71}).

By contrast, Justice Day reasoned, the cotton cloth shipped from the Charlotte mill was ``of themselves, harmless,'' and any harm from child labor had already occurred before interstate transportation began (\href{https://supreme.justia.com/cases/federal/us/247/251/}{247 U.S. at 272}). Relying on \textit{Kidd v. Pearson}, 128 U.S. 1 (1888), \textit{Coe v. Errol}, 116 U.S. 517 (1886), and \textit{Delaware, Lackawanna \& Western R.R. Co. v. Yurkonis}, 238 U.S. 439 (1915), Day drew a rigid categorical line between ``manufacture'' and ``commerce'': production of goods is a purely local matter even when the goods are intended for interstate shipment (\href{https://supreme.justia.com/cases/federal/us/247/251/}{247 U.S. at 272--73}). He also rejected the Government's ``unfair competition'' theory---that states with strict child labor laws were competitively handicapped by states without them---declaring that the Commerce Clause gives Congress no general authority to equalize economic conditions across states (\href{https://supreme.justia.com/cases/federal/us/247/251/}{247 U.S. at 273}). Judging the statute by its ``natural and reasonable effect,'' Day concluded that the Act violated the Constitution in a twofold sense: it exceeded the enumerated commerce power and invaded the local police powers reserved to the States under the Tenth Amendment (\href{https://supreme.justia.com/cases/federal/us/247/251/}{247 U.S. at 274--76}). (Having invalidated the Act on Commerce Clause and Tenth Amendment grounds, the Court did not reach the Fifth Amendment claim.)

\subsection{Concurring Opinions}
No separate concurring opinions were filed.

\subsection{Dissenting Opinion (Justice Holmes)}
Justice Holmes, joined by Justices McKenna, Brandeis, and Clarke, framed the issue starkly: the Keating-Owen Act confined its immediate legal operation to prohibiting the shipment of specified goods in interstate and foreign commerce---an activity squarely within the unqualified text of the Commerce Clause (\href{https://supreme.justia.com/cases/federal/us/247/251/}{247 U.S. at 277--78}). Citing \textit{Champion v. Ames}, Holmes observed that ``regulation means the prohibition of something'' and that the power to regulate commerce includes the power to prohibit part of it outright (\href{https://supreme.justia.com/cases/federal/us/247/251/}{247 U.S. at 277--78}).

Holmes dismantled the majority's indirect-effects objection with a sweeping survey of precedent. In \textit{McCray v. United States} (prohibitive tax on colored oleomargarine), \textit{Veazie Bank v. Fenno} (tax driving state bank notes out of circulation), \textit{Flint v. Stone Tracy Co.} (corporate tax), and Sherman Act antitrust decisions such as \textit{Standard Oil Co. v. United States} and \textit{United States v. American Tobacco Co.}, the Court had repeatedly held that an exercise of an enumerated federal power cannot be invalidated merely because of its obvious indirect effects on state domestic policy or production (\href{https://supreme.justia.com/cases/federal/us/247/251/}{247 U.S. at 278--79}). Holmes rejected the majority's ``harmless goods'' distinction as arbitrary, noting that the Pure Food and Drug Act had been upheld in \textit{Weeks v. United States} even for innocent articles induced by preliminary fraud, and arguing that it does not matter whether the evil precedes or follows interstate transportation so long as Congress concludes that transportation encourages the evil (\href{https://supreme.justia.com/cases/federal/us/247/251/}{247 U.S. at 279--80}). Finally, Holmes emphasized institutional deference: while states remain free to regulate their purely internal affairs, once they seek to send goods across state lines they encounter ``the public policy of the United States, which it is for Congress to express'' (\href{https://supreme.justia.com/cases/federal/us/247/251/}{247 U.S. at 281}).

\section{Examples: Future Applications}
Applying the doctrinal framework announced in \textit{Hammer v. Dagenhart} (as it stood before \href{https://supreme.justia.com/cases/federal/us/312/100/}{\textit{United States v. Darby}} overruled it in 1941) yields the following five applications:

\begin{enumerate}
\item \textbf{Same-Side Example 1 (Unconstitutional Under \textit{Hammer}):} Congress enacts a statute banning the interstate shipment of furniture manufactured in factories where adult workers are paid less than a federal minimum wage or work more than forty hours per week. \textit{Why:} Like the cotton cloth in \textit{Hammer}, the finished furniture is an ordinary, intrinsically harmless commercial commodity, and the labor conditions occur entirely during local manufacture before interstate shipment begins. Under Justice Day's reasoning, Congress may not close the channels of interstate commerce to equalize labor costs across states (\href{https://supreme.justia.com/cases/federal/us/247/251/}{247 U.S. at 272--73}).

\item \textbf{Same-Side Example 2 (Unconstitutional Under \textit{Hammer}):} Congress prohibits the interstate transportation of coal mined in any state that does not require mine operators to contribute to a state-level miners' pension fund. \textit{Why:} Justice Day expressly equated ``the mining of coal'' with ``the making of goods'' as a non-commercial local activity outside Congress's commerce power (\href{https://supreme.justia.com/cases/federal/us/247/251/}{247 U.S. at 272}). Because coal itself is a lawful, ordinary commodity and the regulatory target is local mine employment terms, \textit{Hammer} invalidates the prohibition as an invasion of state police power.

\item \textbf{Opposite-Side Example 1 (Constitutional Under \textit{Hammer}):} Congress prohibits the interstate shipment of toys painted with toxic lead-based paint or meat processed from diseased livestock. \textit{Why:} Unlike child-made cotton cloth, toxic toys and diseased meat are inherently dangerous articles whose interstate shipment directly spreads physical harm into the receiving state. This falls squarely within \textit{Hammer}'s exception for \textit{Hipolite Egg} and other ``noxious goods'' cases where interstate transportation is necessary to accomplish the harmful result (\href{https://supreme.justia.com/cases/federal/us/247/251/}{247 U.S. at 270--71}).

\item \textbf{Opposite-Side Example 2 (Constitutional Under \textit{Hammer}):} Congress prohibits the interstate transportation of stolen automobiles across state lines (as later upheld in \textit{Brooks v. United States}, 267 U.S. 432 (1925)). \textit{Why:} Interstate transportation is actively used as an instrumentality to consummate and conceal a crime, frustrating local law enforcement by removing stolen property beyond state borders. Under \textit{Hammer}'s synthesis of \textit{Champion v. Ames} and \textit{Hoke v. United States}, Congress may bar the use of interstate channels to facilitate illicit schemes (\href{https://supreme.justia.com/cases/federal/us/247/251/}{247 U.S. at 270--71}; \href{https://www.law.cornell.edu/constitution-conan/amendment-10/state-police-power-and-tenth-amendment-jurisprudence}{Cornell LII}).

\item \textbf{Fence-Sitter Example (Borderline Under \textit{Hammer}):} Congress enacts a law banning the interstate shipment of prison-made goods \textit{only} into states whose own laws prohibit the sale or possession of convict-made goods (the scenario later addressed in \href{https://supreme.justia.com/cases/federal/us/299/334/}{\textit{Kentucky Whip \& Collar Co. v. Illinois Central R. Co.}, 299 U.S. 334 (1937)}). \textit{Why:} On the \textit{Hammer} side, prison-made harnesses and garments are physically harmless commodities whose labor conditions occur entirely during pre-shipment manufacture. On the opposite side, the federal ban does not impose a uniform federal labor code on producing states; instead, like the Webb-Kenyon liquor law sustained in \textit{Clark Distilling}, it uses interstate commerce control solely to prevent interstate channels from frustrating the valid local police laws of receiving states---the very distinction the Court later used in \textit{Kentucky Whip} to uphold the statute without yet overruling \textit{Hammer} (\href{https://supreme.justia.com/cases/federal/us/299/334/}{299 U.S. at 346--48}).
\end{enumerate}

\section{Critique}
\textbf{Originalist and Textualist Perspectives:} Originalist and textualist scholars are sharply divided on \textit{Hammer v. Dagenhart}, separating its textual missteps from its structural federalism premise. Textualists universally criticize Justice Day's misquotation of the Tenth Amendment: Day wrote that ``the powers not \textit{expressly} delegated to the National Government are reserved'' to the States (\href{https://supreme.justia.com/cases/federal/us/247/251/}{247 U.S. at 275}), inserting the word ``expressly'' that the Framers deliberately omitted from the Tenth Amendment (unlike Article II of the Articles of Confederation), as Chief Justice Marshall had emphasized in \textit{McCulloch v. Maryland} (\href{https://fedsoc.org/case/hammer-v-dagenhart}{Federalist Society}; \href{https://constitution.congress.gov/browse/essay/amdt10-3-3/ALDE_00013623/}{Constitution Annotated}). At the same time, many modern originalists---including Randy Barnett and Richard Epstein---defend \textit{Hammer}'s core distinction between ``commerce'' (trade, exchange, and transportation) and ``manufacture'' or local production as faithful to Founding-era usage and the doctrine of enumerated powers, arguing that a pretextual ban on interstate shipment aimed solely at regulating local factories cannot be reconciled with limited federal authority (\href{https://digitalcommons.law.byu.edu/cgi/viewcontent.cgi?article=3408&context=lawreview}{BYU Law Review}; \href{https://ilr.law.uiowa.edu/sites/ilr.law.uiowa.edu/files/2024-03/ILR-109-Coan-Schwartz_0.pdf}{Iowa Law Review}). Conversely, broader originalists such as Akhil Reed Amar and Jack Balkin read ``commerce'' as encompassing wider economic and social intercourse, and even some textualists note that Justice Holmes's dissent was textually powerful because the Keating-Owen Act directly regulated interstate shipment itself (\href{https://digitalcommons.law.byu.edu/cgi/viewcontent.cgi?article=3408&context=lawreview}{BYU Law Review}).

\textbf{Progressive and Doctrinal Perspectives:} Progressive legal scholars and doctrinal historians---from early critics such as Thomas Reed Powell and Edward S. Corwin to later constitutional historians---have long condemned \textit{Hammer} as a textbook example of formalistic dual federalism divorced from industrial reality (\href{https://conlaw.jotwell.com/federalism-and-child-labor-revisited/}{Jotwell}; \href{https://www.law.cornell.edu/constitution-conan/article-1/section-8/clause-3/dual-federalism-and-commerce-clause}{Cornell LII}). Doctrinally, critics argue that Day's ``harmful goods'' distinction was incoherent: lottery tickets in \textit{Champion v. Ames} were merely slips of paper, not physically dangerous objects, yet the Court let Congress exclude them because of Congress's moral judgment about gambling, while refusing to let Congress make the same judgment about child labor (\href{https://supreme.justia.com/cases/federal/us/247/251/}{247 U.S. at 280}; \href{https://www.abajournal.com/magazine/article/june_3_1918_child_labor_law_declared_unconstitutional}{ABA Journal}). Economically, progressives emphasize the classic collective-action problem (``race to the bottom''): individual states that banned child labor placed their own textile mills at a competitive disadvantage against states that allowed it, making national regulation the only practical solution---the exact rationale the Court later adopted when unanimously overruling \textit{Hammer} in \href{https://supreme.justia.com/cases/federal/us/312/100/}{\textit{United States v. Darby}, 312 U.S. 100 (1941)}.

\textbf{Underlying Competing Values:} \textit{Hammer v. Dagenhart} exposes a deep structural tension in American constitutionalism between \textit{enumerated-powers federalism} (preserving meaningful state autonomy over local social and economic life so the Commerce Clause does not become a plenary national police power) and \textit{national democratic accountability} (empowering Congress to solve nationwide market-driven collective-action failures and protect vulnerable workers when interstate competition paralyzes individual states).

\section{Key Quotations}
\begin{quote}
``The making of goods and the mining of coal are not commerce, nor does the fact that these things are to be afterwards shipped or used in interstate commerce make their production a part thereof.''
\end{quote}
\noindent --- Justice William R. Day, Opinion of the Court, \href{https://supreme.justia.com/cases/federal/us/247/251/}{247 U.S. at 272} (\href{https://www.courtlistener.com/opinion/99190/hammer-v-dagenhart/}{CourtListener}).

\begin{quote}
``The grant of power to Congress over the subject of interstate commerce was to enable it to regulate such commerce, and not to give it authority to control the States in their exercise of the police power over local trade and manufacture.''
\end{quote}
\noindent --- Justice William R. Day, Opinion of the Court, \href{https://supreme.justia.com/cases/federal/us/247/251/}{247 U.S. at 273--74} (\href{https://www.courtlistener.com/opinion/99190/hammer-v-dagenhart/}{CourtListener}).

\begin{quote}
``In interpreting the Constitution, it must never be forgotten that the Nation is made up of States to which are entrusted the powers of local government. And to them and to the people the powers not expressly delegated to the National Government are reserved.''
\end{quote}
\noindent --- Justice William R. Day, Opinion of the Court, \href{https://supreme.justia.com/cases/federal/us/247/251/}{247 U.S. at 275} (\href{https://www.courtlistener.com/opinion/99190/hammer-v-dagenhart/}{CourtListener}).

\begin{quote}
``But if an act is within the powers specifically conferred upon Congress, it seems to me that it is not made any less constitutional because of the indirect effects that it may have, however obvious it may be that it will have those effects, and that we are not at liberty upon such grounds to hold it void.''
\end{quote}
\noindent --- Justice Oliver Wendell Holmes, Jr., Dissenting, \href{https://supreme.justia.com/cases/federal/us/247/251/}{247 U.S. at 277} (\href{https://www.courtlistener.com/opinion/99190/hammer-v-dagenhart/}{CourtListener}).

\begin{quote}
``It does not matter whether the supposed evil precedes or follows the transportation. It is enough that, in the opinion of Congress, the transportation encourages the evil.''
\end{quote}
\noindent --- Justice Oliver Wendell Holmes, Jr., Dissenting, \href{https://supreme.justia.com/cases/federal/us/247/251/}{247 U.S. at 279--80} (\href{https://www.courtlistener.com/opinion/99190/hammer-v-dagenhart/}{CourtListener}).

\end{document}
